% =============================================================================
% SECTION: Historical Context and Global Rules (English)
% @rpg-schema: <genova1507:> a fitd:Game ;
%   rdfs:label "Genoa 1507 --- The Revolt Against France" ;
%   fitd:basedOn fitd:BladesInTheDark ;
%   rpg:setting "Genoa, 1507" ;
%   rpg:playerCount "16-20" ;
%   rpg:sessionType "one-shot multi-table" ;
%   rpg:sessionDuration "3 hours" .
% =============================================================================

\chapter{Historical Context}
\label{ch:context}

\begin{multicols}{2}

\noindent\textbf{Genoa, 1507.} La Superba --- \textit{the Proud One} ---
has chafed under French occupation for nearly a decade.
Louis~XII holds the city through foreign governors, oppressive taxation,
and garrisons that humiliate its citizens at every turn. But patience has limits.

The great Genoese families --- split by centuries of rivalry between
\idx{Guelphs} and \idx{Ghibellines}, between sea and land --- now face
a decisive choice: who will \textit{lead} the revolt, who will \textit{sabotage}
it for personal gain, and who will sell their soul to France to survive?

Tonight, each table plays one of the four great \idx{Houses}.
The fate of Genoa will be settled before dawn.

\columnbreak

\begin{rulebox}{Historical Note}
The 1507 revolt is well-documented: the Genoese rose against French governor
Gian Giacomo Trivulzio, besieged the castle of Castelletto, and briefly
reclaimed their autonomy. Louis~XII returned militarily in May 1507 and
retook the city. This game explores the narrative space of that night ---
the choices that history compressed into a handful of hours.
\end{rulebox}

\end{multicols}

% =============================================================================
\chapter{Game Structure}
\label{ch:structure}
% =============================================================================

% @rpg-schema: genova1507: rpg:sessionDuration "3 hours" .

\begin{rulebox}[GoldRPG]{Recommended Duration: 3 Hours}
The entire session --- all tables running simultaneously --- is calibrated
for \textbf{approximately 3 hours} of actual play. The 8-segment clock and
two-heist structure (House Heist + Global Heist) are designed to resolve
within this window with 4 active tables.
\end{rulebox}

\section{The Shared Clock}
\label{sec:clock}

% @rpg-schema: genova1507:ControlloGenova a fitd:ProgressClock ;
%   rdfs:label "Control of Genoa" ;
%   fitd:segments "8" ;
%   fitd:shared "true" .

At the centre of every table sits a single \idxbold{Shared Clock}
with 8 segments: \textbf{Control of Genoa}.

\clocktrack[GoldRPG]{CONTROL OF GENOA \ \ [shared clock --- updated by the Coordinator GM]}{0}{8}

\textit{\small The 8-segment clock is paced for \textbf{roughly 3 hours} ---
about 45 minutes per narrative phase (segments 1--2, 3--4, 5--6, 7--8).}

\begin{multicols}{2}

\subsection{Reading the Segments}

\begin{itemize}
  \item[\textbf{1--2}] France holds firm. The houses move in shadow.
  \item[\textbf{3--4}] The revolt smoulders. French officers grow suspicious.
  \item[\textbf{5--6}] Genoa is burning. The \textit{sestieri} choose their champions.
  \item[\textbf{7}]    The French governor flees --- or calls for reinforcements.
  \item[\textbf{8}]    \textbf{GENOA IS FREE} --- but who governs it?
\end{itemize}

\columnbreak

\subsection{The Clock Advances When}
\begin{itemize}
  \item A house completes its \idx{House Heist} successfully \textbf{(+2)}
  \item A table contributes to the \idx{Global Heist} \textbf{(+1 per table)}
  \item An exceptional roll on a narratively significant action \textbf{(+1, GM call)}
\end{itemize}

\subsection{The Clock Falls Back When}
\begin{itemize}
  \item An interference action succeeds \textbf{(--1)}
  \item A heist fails with Serious Consequences \textbf{(--1)}
  \item A house betrays the others to the French \textbf{(--2)}
\end{itemize}

\end{multicols}

% =============================================================================
\section{The Two Heists}
\label{sec:heists}
% =============================================================================

\begin{multicols}{2}

\begin{rulebox}[GoldRPG]{The House Heist}
Every house has a unique House Heist that only \textit{they} can run.
It reflects that family's specific resources, contacts, and history.

\medskip
\textbf{Reward:} +2 to the Shared Clock on success.

\medskip
No coordination with other tables is required --- but another house
can attempt to \textbf{Interfere} (see \S\ref{sec:interference}).
\end{rulebox}

\columnbreak

\begin{rulebox}[GoldRPG]{The Global Heist: The Night of the Genoese Vespers}
Available to all tables simultaneously. Activates when the clock
reaches \textbf{4 segments}.

\medskip
\textbf{Scenario:} Ring the bells of San Lorenzo cathedral at the
agreed signal, triggering a popular uprising across all \textit{sestieri}.

\medskip
Every house plays its role. Count the successes:
\begin{itemize}
  \item \textbf{3+ successes:} +3 to the clock
  \item \textbf{2 successes:} +1 to the clock
  \item \textbf{$\leq$1 success:} --1 to the clock
\end{itemize}

A table may choose to \textbf{Obstruct} instead of cooperate
(see \S\ref{sec:interference}).
\end{rulebox}

\end{multicols}

% =============================================================================
\section{The Interference Rule}
\label{sec:interference}
% =============================================================================

% @rpg-schema: genova1507:Interferenza a fitd:Rule ;
%   rdfs:label "Interference" ;
%   rpg:frequency "once per session per house" .

A house may attempt to disrupt another house's Heist \textbf{once per session}.
The procedure:

\begin{enumerate}
  \item \textbf{Declaration:} The interfering table informs the Coordinator GM,
        who notifies the target table's GM.
  \item \textbf{Interference Roll:} The interfering house rolls
        \textsc{Intrigue} or \textsc{Force} depending on their method.
        The target house may counter with an \textsc{Awareness} roll.
  \item \textbf{Consequences:}
  \begin{itemize}
    \item \textit{Full success:} the target Heist loses 1 clock segment or 1 Resource.
    \item \textit{Partial success:} both houses pay a price.
    \item \textit{Failure:} the interfering house suffers a Complication.
  \end{itemize}
\end{enumerate}

% =============================================================================
\section{House Resources}
\label{sec:resources}
% =============================================================================

\begin{center}
\begin{tabular}{L{2.8cm} C{1.5cm} p{9cm}}
  \toprule
  \textbf{Resource} & \textbf{Scale} & \textbf{Description} \\
  \midrule
  \idx{Coin}        & 0--6 & Ready money for bribes, equipment, and alliances. \\
  \idx{Reputation}  & 0--6 & Political weight in the streets and the palaces. Affects social rolls. \\
  \idx{Influence}   & 0--6 & Access to privileged information and outstanding favours. \\
  \idx{Soldiers}    & 0--6 & Armed men at your disposal. \textit{Consumable.} \\
  \bottomrule
\end{tabular}
\end{center}

% =============================================================================
\section{The Dice System (FitD)}
\label{sec:dice}
% =============================================================================

\begin{multicols}{2}

\subsection{Building the Pool}

Gather d6s equal to the relevant action rating and take
the \textbf{highest} result. In \idx{Desperate Position},
always roll \textbf{1d} regardless of rating.

\begin{center}
\rowcolors{2}{LightBG}{white}
\begin{tabular}{C{1.2cm} p{4.5cm}}
  \toprule
  \textbf{Die} & \textbf{Result} \\
  \midrule
  \textbf{6}   & Full success. No complications. \\
  \textbf{4--5}& Partial success. Complication. \\
  \textbf{1--3}& Failure. Serious complication. \\
  \textbf{66}  & Critical success. Enhanced effect. \\
  \bottomrule
\end{tabular}
\end{center}

\columnbreak

\subsection{Positions}

\begin{description}[leftmargin=2.4cm, labelindent=0pt]
  \item[\idx{Controlled}]  Advantage. Minor consequences.
  \item[\idx{Risky}]       Standard. Moderate consequences.
  \item[\idx{Desperate}]   1d fixed. Serious consequences ---
                           but a success counts double.
\end{description}

\subsection{Resistance}

Spend \textbf{Stress} to resist consequences.
Roll the relevant attribute pool (Insight / Prowess / Resolve):
\begin{itemize}
  \item \textbf{6}: 0 stress \quad \textbf{4--5}: 1 stress \quad \textbf{1--3}: 2 stress
\end{itemize}

\end{multicols}

% =============================================================================
\section{Victory Conditions}
\label{sec:victory}
% =============================================================================

\begin{rulebox}[GoldRPG]{The Clock Reaches 8: GENOA IS FREE}
The house with the most \textbf{Reputation} leads the new government.
Ties broken by \textbf{Influence}. Further ties trigger a short
played epilogue.
\end{rulebox}

\medskip

\begin{center}
\begin{tabular}{L{2.5cm} p{11cm}}
  \toprule
  \textbf{House} & \textbf{Secondary Victory} \\
  \midrule
  \textbf{Doria}   & House Heist complete \textit{and} at least 3 Soldiers remaining:
                     they command the new Republic's navy. \\
  \textbf{Spinola} & Influence 5+ at session end: they control the Republic's bank. \\
  \textbf{Fieschi} & The bells of San Lorenzo rang: they hold ecclesiastical appointments. \\
  \textbf{Grimaldi}& Monaco recognised (House Heist complete): independent victory
                     \textit{regardless of the clock}. \\
  \bottomrule
\end{tabular}
\end{center}

% =============================================================================
\section{A Note on Language}
\label{sec:language}
% =============================================================================

\begin{rulebox}[SharedBrown]{Italian Terms in Play}
Several Italian words are kept untranslated throughout this document because
they carry specific cultural weight that English equivalents would flatten.

\begin{center}
\begin{tabular}{L{3cm} p{9.5cm}}
  \toprule
  \textbf{Term} & \textbf{Meaning in context} \\
  \midrule
  \textit{sestiere}   & City district (Genoa had six: Prè, Portoria, San Giovanni,
                        San Lorenzo, Castello, Soziglia). \\
  \textit{caruggi}    & The narrow medieval alleyways of Genoa's old town ---
                        a labyrinth that only locals navigate freely. \\
  \textit{podestà}    & A city's chief magistrate, often an outsider appointed
                        to ensure impartiality. \\
  \textit{Maona}      & A Genoese commercial consortium, here used for the
                        informal council of merchant interests at the harbour. \\
  \textit{condottiero}& A mercenary captain who commands a hired company. \\
  \bottomrule
\end{tabular}
\end{center}
\end{rulebox}
